%%%%%%%%%%%%%%%%%%%%%%%%%%%%%%%%%%%%%%%%%%%%%%%%%%%%%%%%%%%%%%%%%%%%%%%%%%%
% Función: catálogo de titulaciones soportadas por universidad y centro.
%
% Uso: referencia del manual; normalmente no debes modificarlo.
% Los nombres y traducciones se obtienen del registro de titulaciones.
%%%%%%%%%%%%%%%%%%%%%%%%%%%%%%%%%%%%%%%%%%%%%%%%%%%%%%%%%%%%%%%%%%%%%%%%%%%

\chapter{Titulaciones soportadas}
\label{app:titulaciones-soportadas}

En este apéndice encontrarás los identificadores que puedes asignar a \texttt{\textbackslash{}myDegree}, junto con el nombre completo de cada titulación y, entre paréntesis, su traducción al inglés cuando esté disponible. Si aparece \emph{not translated}, la plantilla no tiene registrado un nombre en inglés para esa titulación. Disponer de una traducción no implica que todas sus páginas institucionales se traduzcan, pues algunas conservan el modelo en español de la universidad. Las variantes de doctorado por compendio se indican expresamente.

Las titulaciones se agrupan por universidad y centro. El soporte de las titulaciones ajenas a la UAH no se ha probado tan exhaustivamente como el de las titulaciones de la UAH. En los grados de la ETSII-URJC y de la EIF-URJC solo se admiten TFG individuales, no colectivos. El informe experimental de GEINTRA no es una titulación y no forma parte de este catálogo.

\DegreeRegistryPrintCatalogue

\section{Indicaciones para los TFG de la EIF-URJC}
\label{sec:indicaciones-eif-urjc}

Además de elegir el identificador de tu titulación, ten en cuenta estas indicaciones al preparar el contenido de tu memoria.

Para los TFG de la EIF-URJC, utiliza el idioma indicado en la guía docente. El reglamento limita la memoria a 100 páginas sin contar anexos y el resumen a una página. En las conclusiones debes mencionar las competencias del grado que has aplicado y las nuevas que has adquirido. También debes incluir una reflexión sobre el impacto social, económico y medioambiental del trabajo y sus responsabilidades éticas y profesionales si tu cohorte de entrada es 2025/2026 o posterior. Para las cohortes anteriores, esta reflexión se recomienda. Los márgenes, la fuente y el interlineado que cita el reglamento son ejemplos de un formato legible, no ajustes impuestos automáticamente por la plantilla. Comprueba siempre los requisitos vigentes con tu tutor o tutora.
